\documentclass[12pt]{article}
\usepackage{sbc-template}
\usepackage{graphicx,url}
\usepackage[utf8]{inputenc}
\usepackage[brazil]{babel}
\usepackage{booktabs}
\usepackage{multirow}
\usepackage{pgfplots}
\pgfplotsset{compat=1.18}

\sloppy

\title{Comparação de Desempenho entre Abordagens Sequencial e Paralela\\
  para Download de Imagens via PokéAPI}

\author{Alexandre Jota Quaresma}

\address{
  Instituto Federal de Pernambuco (IFPE) -- Campus Belo Jardim \\
  Curso de Engenharia de Software \\
  Belo Jardim -- PE -- Brasil
  \email{ajq@discente.ifpe.edu.br}
}

\begin{document}

\maketitle

\begin{abstract}
  This paper presents a performance comparison between sequential and parallel
  approaches for downloading Pokémon images from the PokéAPI REST service.
  Three parallelism strategies were implemented using Python: \texttt{threading},
  \texttt{multiprocessing}, and \texttt{concurrent.futures}. Each approach was
  evaluated with 50, 250, and 500 image requests, repeated 5 times to obtain
  average execution times. Results show that all parallel approaches achieved
  approximately 8x speedup over the sequential baseline when using 8
  workers, with \texttt{threading} and \texttt{concurrent.futures} performing
  nearly identically, and \texttt{multiprocessing} showing slightly higher
  overhead for smaller workloads due to process creation costs.
\end{abstract}

\begin{resumo}
  Este artigo apresenta uma comparação de desempenho entre abordagens sequencial
  e paralela para o download de imagens de Pokémons a partir da API REST PokéAPI.
  Três estratégias de paralelismo foram implementadas em Python: \texttt{threading},
  \texttt{multiprocessing} e \texttt{concurrent.futures}. Cada abordagem foi
  avaliada com 50, 250 e 500 requisições de imagens, repetidas 5 vezes para
  obtenção do tempo médio de execução. Os resultados mostram que todas as
  abordagens paralelas alcançaram aproximadamente 8x de aceleração em relação à
  baseline sequencial com 8 workers, com \texttt{threading} e
  \texttt{concurrent.futures} apresentando desempenho praticamente idêntico, e
  \texttt{multiprocessing} exibindo overhead levemente maior para cargas menores
  devido ao custo de criação de processos.
\end{resumo}


% =============================================================================
\section{Introdução}
% =============================================================================

O avanço da computação paralela e distribuída tornou-se fundamental para o
desenvolvimento de sistemas que realizam operações intensivas de entrada e saída
(I/O), como chamadas a APIs externas e downloads de arquivos. Neste contexto,
comparar o desempenho de diferentes estratégias de paralelismo é essencial para
a tomada de decisões arquiteturais em projetos de software.

Este trabalho tem como objetivo implementar e comparar o desempenho de uma
abordagem sequencial com três abordagens paralelas para o download de imagens de
Pokémons a partir da PokéAPI~\cite{pokeapi}. A PokéAPI é uma API REST gratuita
e de código aberto que disponibiliza dados sobre Pokémons — incluindo habilidades,
índices de jogos, imagens e espécies — servindo mais de 50 bilhões de chamadas
por mês.

As abordagens paralelas foram implementadas utilizando as bibliotecas nativas do
Python: \texttt{threading}, \texttt{multiprocessing} e \texttt{concurrent.futures}.
Para cada abordagem, foram testadas configurações com 2, 4 e 8
\textit{workers}/processos, avaliadas sobre conjuntos de 50, 250 e 500
imagens, com 5 execuções por cenário para obtenção do tempo médio.


% =============================================================================
\section{Fundamentação Teórica}
% =============================================================================

\subsection{PokéAPI}

A PokéAPI disponibiliza o endpoint \url{https://pokeapi.co/api/v2/pokemon/\{id\}}
para recuperar dados de Pokémons por identificador numérico sequencial. A resposta
é retornada em formato JSON e não requer autenticação. Para este trabalho, o campo
de interesse é \texttt{sprites.front\_default}, que contém a URL da imagem
principal do Pokémon, a qual é baixada e armazenada em disco.

\subsection{Paralelismo em Python}

Python oferece três mecanismos principais para execução paralela:

\textbf{threading}: permite a execução concorrente de múltiplas \textit{threads}
dentro de um mesmo processo. Embora o GIL (\textit{Global Interpreter Lock})
limite a paralelização de código Python puro, operações de I/O liberam o GIL,
tornando \textit{threads} eficientes para tarefas de rede~\cite{python-threading}.

\textbf{multiprocessing}: cria processos independentes com espaços de memória
separados, contornando o GIL. É mais adequado para tarefas intensivas em CPU.
Para I/O, o overhead de criação de processos pode reduzir o ganho de
desempenho~\cite{python-multiprocessing}.

\textbf{concurrent.futures}: oferece uma interface de alto nível unificada para
\texttt{ThreadPoolExecutor} e \texttt{ProcessPoolExecutor}, simplificando o
gerenciamento de \textit{pools} de \textit{workers}~\cite{python-futures}.


% =============================================================================
\section{Metodologia}
% =============================================================================

\subsection{Implementação}

Todas as implementações foram desenvolvidas em Python 3.13, utilizando a biblioteca
\texttt{requests} para as requisições HTTP e \texttt{os} para o armazenamento das
imagens em disco. O processo de download de cada Pokémon consiste em duas
requisições: (1) consulta à PokéAPI para obtenção da URL do sprite e (2) download
da imagem a partir dessa URL.

Na abordagem \textbf{sequencial}, os downloads são realizados um por vez, do
Pokémon de id 1 ao id $N$. Nas abordagens \textbf{paralelas}, o número de
\textit{workers} simultâneos é controlado — por \texttt{Semaphore} no
\texttt{threading}, por \texttt{Pool} no \texttt{multiprocessing} e por
\texttt{ThreadPoolExecutor} no \texttt{concurrent.futures}.

\subsection{Protocolo de Avaliação}

Para garantir a imparcialidade das medições, o diretório de imagens é removido e
recriado antes de cada rodada, forçando o download completo a cada execução e
eliminando a influência de cache em disco. O tempo é medido com
\texttt{time.time()}, capturando o intervalo desde o início até o término de
todos os downloads.

Cada cenário foi executado 5 vezes. A métrica reportada é o \textbf{tempo médio
de execução} em segundos.

Os testes foram executados na seguinte configuração de hardware e software:

\begin{itemize}
  \item Sistema Operacional: Windows 11
  \item Processador: % TODO: preencher (ex: Intel Core i5-XXXX, X núcleos)
  \item Memória RAM: % TODO: preencher (ex: 16 GB)
  \item Python 3.13
  \item Conexão com a Internet: % TODO: preencher (ex: 100 Mbps)
\end{itemize}


% =============================================================================
\section{Resultados}
% =============================================================================

As Tabelas~\ref{tab:seq}, \ref{tab:thread}, \ref{tab:multiproc}
e~\ref{tab:futures} apresentam os tempos médios de execução (em segundos) para
cada abordagem avaliada, seguidas dos gráficos comparativos.

% --- Tabela 1: Sequencial ---
\begin{table}[ht]
\centering
\caption{Tempo médio de execução — Abordagem Sequencial (5 rodadas)}
\label{tab:seq}
\begin{tabular}{cc}
\toprule
\textbf{Pokémons} & \textbf{Tempo Médio (s)} \\
\midrule
50  & 9.27 \\
250 & 45.35 \\
500 & 93.82 \\
\bottomrule
\end{tabular}
\end{table}

% --- Tabela 2: Threading ---
\begin{table}[ht]
\centering
\caption{Tempo médio de execução — \texttt{threading} (5 rodadas)}
\label{tab:thread}
\begin{tabular}{cccc}
\toprule
\textbf{Pokémons} & \textbf{2 Threads (s)} & \textbf{4 Threads (s)} & \textbf{8 Threads (s)} \\
\midrule
50  & 4.15 & 1.99 & 1.14 \\
250 & 22.70 & 10.36 & 5.57 \\
500 & 46.39 & 21.92 & 12.02 \\
\bottomrule
\end{tabular}
\end{table}

% --- Tabela 3: Multiprocessing ---
\begin{table}[ht]
\centering
\caption{Tempo médio de execução — \texttt{multiprocessing} (5 rodadas)}
\label{tab:multiproc}
\begin{tabular}{cccc}
\toprule
\textbf{Pokémons} & \textbf{2 Processos (s)} & \textbf{4 Processos (s)} & \textbf{8 Processos (s)} \\
\midrule
50  & 5.05 & 2.47 & 1.71 \\
250 & 23.84 & 11.45 & 5.82 \\
500 & 49.25 & 22.89 & 11.86 \\
\bottomrule
\end{tabular}
\end{table}

% --- Tabela 4: Concurrent.futures ---
\begin{table}[ht]
\centering
\caption{Tempo médio de execução — \texttt{concurrent.futures} (5 rodadas)}
\label{tab:futures}
\begin{tabular}{cccc}
\toprule
\textbf{Pokémons} & \textbf{2 Workers (s)} & \textbf{4 Workers (s)} & \textbf{8 Workers (s)} \\
\midrule
50  & 4.16 & 1.99 & 1.13 \\
250 & 21.88 & 10.55 & 5.64 \\
500 & 47.07 & 22.40 & 12.03 \\
\bottomrule
\end{tabular}
\end{table}

% --- Gráfico 1: Sequencial vs Threading ---
\begin{figure}[ht]
\centering
\begin{tikzpicture}
\begin{axis}[
    title={Sequencial vs \texttt{threading}},
    xlabel={Número de Pokémons},
    ylabel={Tempo médio (s)},
    xtick={50, 250, 500},
    legend pos=north west,
    ymajorgrids=true,
    grid style=dashed,
    width=0.85\textwidth,
    height=6cm,
]
\addplot[color=black, mark=square, thick]
    coordinates {(50,9.27)(250,45.35)(500,93.82)};
\addplot[color=blue, mark=o, thick]
    coordinates {(50,4.15)(250,22.70)(500,46.39)};
\addplot[color=red, mark=triangle, thick]
    coordinates {(50,1.99)(250,10.36)(500,21.92)};
\addplot[color=teal, mark=diamond, thick]
    coordinates {(50,1.14)(250,5.57)(500,12.02)};
\legend{Sequencial, 2 threads, 4 threads, 8 threads}
\end{axis}
\end{tikzpicture}
\caption{Comparação de tempo médio entre a abordagem sequencial e \texttt{threading}}
\label{fig:seq_vs_thread}
\end{figure}

% --- Gráfico 2: Sequencial vs Multiprocessing ---
\begin{figure}[ht]
\centering
\begin{tikzpicture}
\begin{axis}[
    title={Sequencial vs \texttt{multiprocessing}},
    xlabel={Número de Pokémons},
    ylabel={Tempo médio (s)},
    xtick={50, 250, 500},
    legend pos=north west,
    ymajorgrids=true,
    grid style=dashed,
    width=0.85\textwidth,
    height=6cm,
]
\addplot[color=black, mark=square, thick]
    coordinates {(50,9.27)(250,45.35)(500,93.82)};
\addplot[color=blue, mark=o, thick]
    coordinates {(50,5.05)(250,23.84)(500,49.25)};
\addplot[color=red, mark=triangle, thick]
    coordinates {(50,2.47)(250,11.45)(500,22.89)};
\addplot[color=teal, mark=diamond, thick]
    coordinates {(50,1.71)(250,5.82)(500,11.86)};
\legend{Sequencial, 2 processos, 4 processos, 8 processos}
\end{axis}
\end{tikzpicture}
\caption{Comparação de tempo médio entre a abordagem sequencial e \texttt{multiprocessing}}
\label{fig:seq_vs_multi}
\end{figure}

% --- Gráfico 3: Sequencial vs Concurrent.futures ---
\begin{figure}[ht]
\centering
\begin{tikzpicture}
\begin{axis}[
    title={Sequencial vs \texttt{concurrent.futures}},
    xlabel={Número de Pokémons},
    ylabel={Tempo médio (s)},
    xtick={50, 250, 500},
    legend pos=north west,
    ymajorgrids=true,
    grid style=dashed,
    width=0.85\textwidth,
    height=6cm,
]
\addplot[color=black, mark=square, thick]
    coordinates {(50,9.27)(250,45.35)(500,93.82)};
\addplot[color=blue, mark=o, thick]
    coordinates {(50,4.16)(250,21.88)(500,47.07)};
\addplot[color=red, mark=triangle, thick]
    coordinates {(50,1.99)(250,10.55)(500,22.40)};
\addplot[color=teal, mark=diamond, thick]
    coordinates {(50,1.13)(250,5.64)(500,12.03)};
\legend{Sequencial, 2 workers, 4 workers, 8 workers}
\end{axis}
\end{tikzpicture}
\caption{Comparação de tempo médio entre a abordagem sequencial e \texttt{concurrent.futures}}
\label{fig:seq_vs_futures}
\end{figure}

% --- Gráfico 4: Melhor configuração de cada abordagem (8 workers) ---
\begin{figure}[ht]
\centering
\begin{tikzpicture}
\begin{axis}[
    title={Comparação geral — melhor configuração (8 workers/processos)},
    xlabel={Número de Pokémons},
    ylabel={Tempo médio (s)},
    xtick={50, 250, 500},
    legend pos=north west,
    ymajorgrids=true,
    grid style=dashed,
    width=0.85\textwidth,
    height=6cm,
]
\addplot[color=black, mark=square, thick]
    coordinates {(50,9.27)(250,45.35)(500,93.82)};
\addplot[color=blue, mark=o, thick]
    coordinates {(50,1.14)(250,5.57)(500,12.02)};
\addplot[color=red, mark=triangle, thick]
    coordinates {(50,1.71)(250,5.82)(500,11.86)};
\addplot[color=teal, mark=diamond, thick]
    coordinates {(50,1.13)(250,5.64)(500,12.03)};
\legend{Sequencial, Threading (8), Multiprocessing (8), Futures (8)}
\end{axis}
\end{tikzpicture}
\caption{Comparação geral entre todas as abordagens na melhor configuração (8 workers/processos)}
\label{fig:comparacao_geral}
\end{figure}


% =============================================================================
\section{Discussão}
% =============================================================================

Os resultados obtidos confirmam que o paralelismo traz ganhos expressivos para
tarefas de I/O intensivo, como o download de imagens via API REST. A abordagem
sequencial levou 93,82 segundos para baixar 500 imagens, enquanto todas as
versões paralelas com 8 \textit{workers} concluíram a mesma tarefa em torno de
12 segundos — uma aceleração de aproximadamente \textbf{7,8 vezes}.

As abordagens \texttt{threading} e \texttt{concurrent.futures} apresentaram
desempenho praticamente idêntico em todos os cenários. Isso era esperado, pois o
\texttt{concurrent.futures} com \texttt{ThreadPoolExecutor} utiliza
\textit{threads} internamente, sendo essencialmente uma interface de mais alto
nível para o mesmo mecanismo. Ambas se beneficiam diretamente da liberação do
GIL durante as operações de rede, permitindo que múltiplas \textit{threads}
realizem requisições HTTP simultaneamente sem bloqueio mútuo.

O \texttt{multiprocessing} apresentou desempenho competitivo, porém com overhead
levemente maior em cargas menores. Para 50 imagens, o tempo com 2 processos
(5,05s) foi superior ao de 2 \textit{threads} (4,15s), evidenciando o custo de
inicialização dos processos filhos. À medida que a carga aumenta, esse overhead
se torna proporcionalmente menor — com 500 imagens e 8 workers, o
\texttt{multiprocessing} (11,86s) ficou marginalmente mais rápido que o
\texttt{threading} (12,02s), provavelmente devido a uma melhor distribuição de
requisições entre processos com espaços de memória isolados.

Quanto ao número de \textit{workers}, observou-se escalabilidade quase linear
ao dobrar de 2 para 4 e de 4 para 8. Para 500 imagens com \texttt{threading},
o tempo caiu de 46,39s (2 threads) para 21,92s (4 threads) e para 12,02s (8
threads), mantendo uma proporção próxima de 2x a cada duplicação. Isso indica
que, neste cenário, o gargalo ainda é o número de requisições concorrentes, e
não a largura de banda ou o \textit{rate limiting} da API.


% =============================================================================
\section{Conclusão}
% =============================================================================

Este trabalho comparou o desempenho de abordagens sequencial e paralelas para o
download de imagens via PokéAPI, utilizando as bibliotecas \texttt{threading},
\texttt{multiprocessing} e \texttt{concurrent.futures} do Python. Os experimentos
demonstraram que todas as abordagens paralelas reduziram o tempo de execução em
aproximadamente 8 vezes em relação à versão sequencial, com 8 \textit{workers}.

Para cargas de trabalho com gargalo em I/O de rede, \texttt{threading} e
\texttt{concurrent.futures} provaram ser as escolhas mais eficientes e de menor
complexidade de implementação. O \texttt{multiprocessing}, embora competitivo em
cargas maiores, apresenta overhead de inicialização que o torna menos vantajoso
para tarefas menores.

Como trabalho futuro, sugere-se a avaliação de abordagens assíncronas com
\texttt{asyncio} e \texttt{aiohttp}, que podem oferecer desempenho superior para
I/O intensivo sem o overhead de múltiplas \textit{threads} ou processos,
utilizando um único \textit{event loop} para gerenciar centenas de requisições
concorrentes.


% =============================================================================
\bibliographystyle{sbc}
\begin{thebibliography}{9}

\bibitem{pokeapi}
  Hallett, P. et al.
  \textbf{PokéAPI — The RESTful Pokémon API}.
  Disponível em: \url{https://pokeapi.co}.
  Acesso em: maio de 2026.

\bibitem{python-threading}
  Python Software Foundation.
  \textbf{threading — Thread-based parallelism}.
  Disponível em: \url{https://docs.python.org/3/library/threading.html}.
  Acesso em: maio de 2026.

\bibitem{python-multiprocessing}
  Python Software Foundation.
  \textbf{multiprocessing — Process-based parallelism}.
  Disponível em: \url{https://docs.python.org/3/library/multiprocessing.html}.
  Acesso em: maio de 2026.

\bibitem{python-futures}
  Python Software Foundation.
  \textbf{concurrent.futures — Launching parallel tasks}.
  Disponível em: \url{https://docs.python.org/3/library/concurrent.futures.html}.
  Acesso em: maio de 2026.

\end{thebibliography}

\end{document}
